\vspace*{1.5cm}
\section{Resumen}

Se presenta un resumen del Trabajo Práctico 1, detallando los temas a tratar para ofrecer una lectura fluida, clara y precisa que permita comprender el desarrollo en su totalidad.

En primer lugar, se va a hablar sobre el problema propuesto, cómo fue el planteo de la resolución y cómo se logró llegar a que ésta misma pueda ser óptima para cualquier caso posible.

Se detalla de manera concisa la teoría necesaria sobre los \textit{algoritmos Greedy}, enfocándose en lo más básico de la teoría, entre ella la \textit{Regla Greedy}, la sucesión de los óptimos locales, el estado actual y el óptimo global. Para ilustrar estos aspectos, se incluye un breve ejemplo de la idea en general y su analogía a conceptos geopolíticos (como el deontologismo).

Por otro lado, se detalla el diseño del algoritmo implementado en lenguaje \textit{Python}, anteriormente explicando bien la lógica, cómo fue buscada y pensada esa resolución. Adjuntándose al final la implementación realizada.

Se realiza un \textit{análisis} exhaustivo del algoritmo planteado, mediante el estudio de la complejidad y una demostración matemática para validar la optimalidad del algoritmo, de forma clara y con ejemplos propuestos. Se incorporan múltiples experimentos empíricos de forma tal que se pueda visualizar gráficamente las medidas de tiempo sobre cierta cantidad de casos.

Finalmente, se expone una conclusión sobre lo aprendido, la dificultad del problema, y el trayecto.

\textbf{ITEMS PARA HACER:}
\begin{itemize}
    \item  Hacer un análisis del problema, y proponer un algoritmo greedy que obtenga la solución óptima al problema planteado: Dados los n valores de todas las monedas, indicar qué monedas debe ir eligiendo Sophia para si misma y para Mateo, de tal forma que se asegure de ganar siempre. Considerar que Sophia siempre comienza (para sí misma).
    \item  Demostrar que el algoritmo planteado obtiene siempre la solución óptima (desestimando el caso de una cantidad par de monedas de mismo valor, en cuyo caso siempre sería empate más allá de la estrategia de Sophia).
    \item  Escribir el algoritmo planteado. Describir y justificar la complejidad de dicho algoritmo. Analizar si (y cómo) afecta la variabilidad de los valores de las diferentes monedas a los tiempos del algoritmo planteado.
    \item  Analizar si (y cómo) afecta la variabilidad de los valores de las diferentes monedas a la optimalidad del algoritmo planteado.
    \item  Realizar ejemplos de ejecución para encontrar soluciones y corroborar lo encontrado. Adicionalmente, el curso proveerá con algunos casos particulares que deben cumplirse su optimalidad también.
    \item  Hacer mediciones de tiempos para corroborar la complejidad teórica indicada. Agregar los casos de prueba necesarios para dicha corroboración. Esta corroboración empírica debe realizarse confeccionando gráficos correspondientes, y utilizando la técnica de cuadrados mínimos. Para esto, proveemos una explicación detallada, en conjunto de ejemplos.
    \item  Agregar cualquier conclusión que les parezca relevante.
\end{itemize}

